\documentclass[12pt,a4paper]{article}

\usepackage[utf8]{inputenc}
\usepackage[T1]{fontenc}
\usepackage[brazil]{babel}
\usepackage{amsmath,amssymb,amsthm}
\usepackage[a4paper,margin=2.5cm]{geometry}

\newcommand{\R}{\mathbb{R}}
\newcommand{\N}{\mathbb{N}}
\newcommand{\Lp}[2]{L^{#1,#2}}
\newcommand{\norm}[2]{\left\|#1\right\|_{#2}}

\title{Analise Harmonica: exercicio 3.4 da lista 3}
\author{}
\date{}

\begin{document}

\maketitle

Este arquivo e uma referencia textual local para o exercicio 3.4 de
\texttt{lista3.pdf}, usando a notacao de \texttt{Notas\_Harmonica\_2026\_2.pdf}.
Assim, a solucao pode ser consultada sem reabrir os PDFs.

\section*{Enunciado}

Sejam $p_0<p<p_1$, $1\leq q\leq\infty$ e $f\in L^{p,q}$.

\begin{enumerate}
\item Provar que $f=f_L+f_H$, com
\[
\norm{f_L}{L^{p_0,\infty}}+\norm{f_H}{L^{p_1,\infty}}\lesssim \norm{f}{L^{p,q}}.
\]
\item Usar o item anterior e o lema de Fatou para mostrar que, se $T$ e do
tipo fraco $(p_0,q_0)$ e $(p_1,q_1)$, entao $T$ e enumeravelmente subaditivo
na classe de Lorentz intermediaria. Completar a passagem de funcoes simples
para uma funcao geral na prova do teorema de Hunt, sem supor linearidade.
\item Provar que a estimativa restrita
\[
\norm{T\chi_E}{L^{q,\infty}}\lesssim \mu(E)^{1/p}
\]
para conjuntos $E$ de medida finita, nos pares extremos, e equivalente a
\[
\norm{Tf}{L^{q,\infty}}\lesssim \norm{f}{L^{p,1}}.
\]
\end{enumerate}

As constantes abaixo dependem apenas dos expoentes; escrevemos $\lesssim$ para
ocultar tais constantes.

\section*{Notacao das notas}

Para $p<\infty$,
\[
\norm{f}{L^{p,q}}_* = p^{1/q}
\left\|\lambda\,\mu(\{|f|\geq\lambda\})^{1/p}\right\|_{L^q((0,\infty),d\lambda/\lambda)},
\]
e
\[
\norm{f}{L^{p,\infty}}_* =
\sup_{\lambda>0}\lambda\,\mu(\{|f|\geq\lambda\})^{1/p}.
\]

As quase-normas com e sem estrela sao equivalentes. O Lema 3.12 das notas
fornece uma decomposicao diadica: existem $H_n$ decrescentes e conjuntos
disjuntos $E_n$, com $\mu(E_n)\leq 2^n$, tais que, a menos de modificar os
valores nos conjuntos de nivel,
\[
f=\sum_n f_n, \qquad |f_n|\leq H_n\chi_{E_n},
\]
e
\[
a_n=H_n2^{n/p}, \qquad
\|a\|_{\ell^q}\lesssim \norm{f}{L^{p,q}}_*.
\]

O caso $q=\infty$ significa
\[
\sup_n|a_n|\lesssim \norm{f}{L^{p,\infty}}_*.
\]

\section*{1. Decomposicao em dois extremos fracos}

Defina
\[
f_L=\sum_{n\leq 0}f_n, \qquad f_H=\sum_{n>0}f_n.
\]

Primeiro normalize $\norm{f}{L^{p,q}}_*=1$; o caso geral segue multiplicando
a decomposicao por essa norma. Como $|a_n|\leq\|a\|_{\ell^q}\leq 1$, temos
\[
H_n\leq 2^{-n/p}.
\]

\subsection*{Estimativa de $f_L$}

Se $0<\lambda\leq H_0$, entao
\[
\mu(\{|f_L|>\lambda\})
\leq\sum_{n\leq 0}\mu(E_n)
\leq\sum_{n\leq 0}2^n\leq 2,
\]
logo
\[
\lambda\,\mu(\{|f_L|>\lambda\})^{1/p_0}\lesssim 1.
\]

Se $\lambda>H_0$, seja $m\leq 0$ o ultimo indice com $H_m>\lambda$. Entao
\[
\mu(\{|f_L|>\lambda\})
\leq\sum_{n\leq m}2^n\leq 2^{m+1}
\]
e, como $\lambda<H_m\leq 2^{-m/p}$, temos $2^m\leq\lambda^{-p}$. Portanto,
\[
\lambda\,\mu(\{|f_L|>\lambda\})^{1/p_0}
\lesssim \lambda^{1-p/p_0}\leq 1,
\]
pois $p_0<p$ e $\lambda>H_0$, com $H_0\leq 1$.

Assim,
\[
\norm{f_L}{L^{p_0,\infty}}_*\lesssim 1.
\]

\subsection*{Estimativa de $f_H$}

Se $\lambda>H_1$, o conjunto de nivel e vazio. Se $0<\lambda\leq H_1$,
seja $m\geq 1$ o ultimo indice com $H_m>\lambda$. Como antes,
\[
\mu(\{|f_H|>\lambda\})
\leq\sum_{1\leq n\leq m}2^n\leq 2^{m+1}
\]
e $2^m\leq\lambda^{-p}$. Logo
\[
\lambda\,\mu(\{|f_H|>\lambda\})^{1/p_1}
\lesssim \lambda^{1-p/p_1}\leq 1,
\]
porque $p<p_1$ e $\lambda\leq H_1\leq 1$.

Concluimos que
\[
\norm{f_H}{L^{p_1,\infty}}_*\lesssim 1.
\]
Desfazendo a normalizacao,
\[
\norm{f_L}{L^{p_0,\infty}}+\norm{f_H}{L^{p_1,\infty}}\lesssim \norm{f}{L^{p,q}}.
\]

O ponto essencial e separar os blocos diadicos pelo indice $n=0$, e nao por
um corte fixo nos valores de $|f|$. Um corte fixo nao controla a cauda de baixa
amplitude em um espaco de medida infinito.

\section*{2. Subaditividade enumeravel e passagem para funcoes gerais}

Suponha que $T$ seja sublinear:
\[
|T(u+v)|\leq |T(u)|+|T(v)|, \qquad |T(cu)|=|c|\,|T(u)|.
\]

Para uma soma finita $F_N=\sum_{k=1}^Nf_k$, aplique o item 1 a cada parcela:
\[
f_k=f_{k,L}+f_{k,H}.
\]
Entao
\[
|T(F_N)|\leq\sum_{k=1}^N|T(f_{k,L})|
+\sum_{k=1}^N|T(f_{k,H})|.
\]

As hipoteses de tipo fraco nos extremos dao
\[
\norm{T(f_{k,L})}{L^{q_0,\infty}}
\lesssim \norm{f_{k,L}}{L^{p_0,\infty}},
\]
\[
\norm{T(f_{k,H})}{L^{q_1,\infty}}
\lesssim \norm{f_{k,H}}{L^{p_1,\infty}}.
\]

Aplicando a quase-desigualdade triangular de Lorentz as somas finitas e depois
o item 1, obtem-se a estimativa necessaria para $F_N$. Em particular, a
constante nao depende de $N$.

Para a soma infinita, ponha $F_N=\sum_{k=1}^Nf_k$. A sublinearidade fornece,
para cada soma finita,
\[
|T(F_N)|\leq G_N, \qquad G_N=\sum_{k=1}^N|T(f_k)|.
\]

As funcoes $G_N$ crescem pontualmente para
$G=\sum_{k=1}^{\infty}|T(f_k)|$. O Lema de Fatou aplicado a
\[
\lambda\longmapsto
\lambda\,\mu(\{|T(F_N)|>\lambda\})^{1/q}
\]
permite passar ao limite, depois de usar a definicao de $T(\sum f_k)$ como o
limite apropriado das somas finitas (ou a desigualdade sublinear diretamente
com o resto). Assim,
\[
\norm{T(\sum_{k=1}^{\infty}f_k)}{L^{q,\infty}}
\lesssim\sum_{k=1}^{\infty}\norm{T(f_k)}{L^{q,\infty}},
\]
com a interpretacao de quase-norma e a constante da quase-triangularidade de
Lorentz. Este e o detalhe que evita aplicar diretamente a desigualdade
triangular a uma soma infinita.

\subsection*{Extensao de simples para $f\in L^{p,q}$}

Escolha funcoes simples $s_j$ tais que
$\norm{f-s_j}{L^{p,q}}\to 0$ (por truncacao dos valores e do suporte). Escreva
\[
f=s_j+(f-s_j).
\]

Sublinearidade fornece
\[
|Tf|\leq |T(s_j)|+|T(f-s_j)|.
\]

O argumento ja provado para funcoes simples controla o primeiro termo, e a
mesma estimativa aplicada ao resto controla o segundo por
$C\norm{f-s_j}{L^{p,q}}$. Fatou, ou primeiro a passagem para somas finitas e
depois Fatou, da
\[
\norm{Tf}{L^{q,\infty}}\leq C\liminf_{j\to\infty}\norm{s_j}{L^{p,q}}
=C\norm{f}{L^{p,q}}.
\]

Nenhuma identidade $T(s_j)-T(s_k)=T(s_j-s_k)$ foi usada; isso seria ilegitimo
sem linearidade. Usa-se somente a desigualdade sublinear para decompor $f$ em
aproximante mais resto.

\section*{3. Tipo fraco restrito equivale a $L^{p,1}\to L^{q,\infty}$}

\subsection*{Implicacao facil}

Se
\[
\norm{T\chi_E}{L^{q,\infty}}\lesssim\mu(E)^{1/p}
\]
para todo conjunto mensuravel $E$ de medida finita, tome $f=\chi_E$. Como
\[
\norm{\chi_E}{L^{p,1}}=\mu(E)^{1/p},
\]
segue imediatamente que
\[
\norm{Tf}{L^{q,\infty}}\lesssim\norm{f}{L^{p,1}}
\]
na classe das funcoes caracteristicas.

\subsection*{Implicacao geral}

Comece com uma funcao simples nao-negativa, escrita com niveis diadicos:
\[
f=\sum_ka_k\chi_{E_k},
\]
onde os $E_k$ sao disjuntos, $\mu(E_k)<\infty$ e $a_k>0$. Pela sublinearidade
e homogeneidade,
\[
|Tf|\leq\sum_ka_k|T\chi_{E_k}|.
\]

Para uma soma finita, a quase-desigualdade triangular em $L^{q,\infty}$ e a
hipotese restrita fornecem
\[
\norm{T(\sum_{k=1}^Na_k\chi_{E_k})}{L^{q,\infty}}
\lesssim\sum_{k=1}^Na_k\mu(E_k)^{1/p}.
\]

Agora aplique Fatou as somas parciais. Nao se deve escrever diretamente uma
desigualdade triangular para a soma infinita: primeiro controla-se cada soma
finita por uma constante independente de $N$, e somente entao passa-se ao
limite. Pela caracterizacao atomica de Lorentz,
\[
\norm{f}{L^{p,1}}\lesssim\sum_ka_k\mu(E_k)^{1/p},
\]
e portanto
\[
\norm{Tf}{L^{q,\infty}}\lesssim\norm{f}{L^{p,1}}.
\]

Para uma funcao geral $f\in L^{p,1}$, escolha simples nao-negativas $f_j$ que
convergem a $f$ em $L^{p,1}$ e repita o argumento com $f_j$ e o resto $f-f_j$. 
Fatou conclui a estimativa para $f$. Para funcao complexa, aplique o argumento
a $|f|$ quando $|Tf|\leq T(|f|)$ (a forma usual da sublinearidade); caso
contrario, use a decomposicao em partes real e imaginaria prevista na hipotese
sobre $T$.

A necessidade e a suficiencia estao provadas:
\[
T\text{ e do tipo fraco restrito }(p,q)\text{ em conjuntos de medida finita}
\]
se e somente se
\[
\norm{Tf}{L^{q,\infty}}\lesssim\norm{f}{L^{p,1}}.
\]

\section*{Observacao sobre o teorema de Hunt}

Na prova do teorema de Hunt, os pares extremos sao $(p_0,q_0)$ e $(p_1,q_1)$,
e os expoentes intermediarios satisfazem
\[
\frac1p=\frac{1-\theta}{p_0}+\frac{\theta}{p_1},
\qquad
\frac1q=\frac{1-\theta}{q_0}+\frac{\theta}{q_1}.
\]

O item 1 fornece a decomposicao nos dois extremos; o item 2 permite tratar
somas e limites sem linearidade; e o item 3 transforma as estimativas em
conjuntos caracteristicos na estimativa completa $L^{p,1}\to L^{q,\infty}$.
Esses sao precisamente os detalhes omitidos que a lista pede completar.

\end{document}